\chapter{Dữ Liệu}
\label{chap:data}

Dữ liệu trong đồ án gồm hai phần độc lập phục vụ hai module khác nhau của hệ thống:
(1)~Dataset cử chỉ -- cảm xúc tự thu thập, dùng cho Gesture Emotion Encoder và Conditioned Decoder;
(2)~Lakh MIDI Clean Dataset, dùng cho Music Prior.

% ────────────────────────────────────────────────────────────────────
\section{Dataset Cử Chỉ -- Cảm Xúc (Tự Thu Thập)}
% ────────────────────────────────────────────────────────────────────

\subsection{Lý Do Tự Thu Thập}

Không có dataset chuẩn công khai nào chứa cặp dữ liệu
(chuỗi landmarks cơ thể -- nhãn cảm xúc âm nhạc).
Các dataset cảm xúc phổ biến (DEAP, MAHNOB-HCI, IEMOCAP) sử dụng
tín hiệu sinh lý (EEG, GSR, EMG) hoặc dữ liệu giọng nói,
không phù hợp với bài toán điều khiển âm nhạc bằng cử chỉ.
Do đó, nhóm thiết kế và thu thập dataset riêng theo giao thức có kiểm soát.

\subsection{Không Gian Nhãn Cảm Xúc}

Nhóm sử dụng mô hình \textbf{Valence-Arousal}~\cite{russell1980} ---
khung biểu diễn cảm xúc 2 chiều được sử dụng rộng rãi trong nghiên cứu
âm nhạc tình cảm (Music Emotion Recognition):
\begin{align}
  v &\in [-1, +1] \quad \text{(Valence: buồn} \leftrightarrow \text{vui)} \\
  a &\in [-1, +1] \quad \text{(Arousal: thư giãn} \leftrightarrow \text{kích động)}
\end{align}

Mô hình này phù hợp với đồ án vì:
(i)~có thể ánh xạ trực tiếp sang đặc tính âm nhạc (tempo, scale, dynamics);
(ii)~người dùng có thể biểu đạt tự nhiên qua chuyển động cơ thể.

\begin{figure}[H]
  \centering
  \includegraphics[width=0.9\textwidth]{images/data/fig01_emotion_space.png}
  \caption{Phân bố toàn bộ dataset trong không gian Valence-Arousal.
           (a)~Scatter plot tất cả mẫu, phân biệt màu theo chế độ cảm xúc.
           (b)~Trung tâm và độ lệch chuẩn của mỗi chế độ cho thấy
           các chế độ phân tách rõ ràng trong không gian 2 chiều.}
  \label{fig:emotion_space}
\end{figure}

\subsection{Thiết Kế 6 Chế Độ Thu Thập}

Nhóm chia không gian Valence-Arousal thành 6 vùng đặc trưng,
đảm bảo phủ đều các góc phần tư và điểm gốc:

\begin{table}[H]
  \centering
  \caption{Sáu chế độ cảm xúc, nhãn Valence-Arousal và đặc trưng cử chỉ}
  \label{tab:modes}
  \begin{tabular}{clccl}
    \toprule
    \textbf{ID} & \textbf{Tên chế độ} & \textbf{Valence} & \textbf{Arousal}
                & \textbf{Đặc trưng cử chỉ} \\
    \midrule
    0 & HAPPY\_HIGH  & $+1.0$ & $+1.0$
      & Tay nhanh, rộng, lên cao (như cổ vũ) \\
    1 & HAPPY\_LOW   & $+1.0$ & $-0.5$
      & Tay chậm, chuyển động tròn, nhẹ nhàng \\
    2 & SAD\_HIGH    & $-1.0$ & $+0.8$
      & Tay thấp, giật cục, biên độ nhỏ \\
    3 & SAD\_LOW     & $-1.0$ & $-1.0$
      & Tay ở thấp nhất, rất chậm, nặng nề \\
    4 & NEUTRAL      &  $0.0$ &  $0.0$
      & Tay ở giữa màn hình, di chuyển đều \\
    5 & NONE         &  $0.0$ &  $0.0$
      & Không có tay trong khung hình \\
    \bottomrule
  \end{tabular}
\end{table}

\subsection{Đặc Trưng Đầu Vào (Input Features)}

MediaPipe Holistic~\cite{mediapipe2019} trích xuất ba nhóm điểm mốc:
\begin{itemize}
  \item 21 điểm bàn tay phải $\times$ 3 tọa độ $(x, y, z)$ = 63 features
  \item 21 điểm bàn tay trái $\times$ 3 tọa độ = 63 features
  \item 33 điểm toàn thân (pose landmarks) $\times$ 3 tọa độ = 99 features
\end{itemize}
Tổng: \textbf{225 features mỗi frame}.

Tất cả được normalize về tọa độ tương đối:
\begin{equation}
  \tilde{\mathbf{p}}_i = \frac{\mathbf{p}_i - \mathbf{p}_{\text{ref}}}{\|\mathbf{p}_{\max} - \mathbf{p}_{\min}\|_2 + \epsilon}
\end{equation}
trong đó $\mathbf{p}_{\text{ref}}$ là điểm cổ tay/hông, $\epsilon = 10^{-6}$ để tránh chia 0.
Cách normalize này đảm bảo tính bất biến với khoảng cách camera và tầm vóc người dùng.

\subsection{Giao Thức Thu Thập}

Mỗi người thu thập trải qua quy trình sau cho từng chế độ:
\begin{enumerate}
  \item Xem màn hình giới thiệu chế độ, hướng dẫn cử chỉ bằng hình ảnh minh họa.
  \item Nhấn \texttt{SPACE} --- hệ thống đếm ngược 3-2-1.
  \item \textbf{Nghe nhạc nền} (được thiết kế theo cảm xúc của chế độ) và
        di chuyển tay tự nhiên theo cảm xúc của nhạc.
  \item Hệ thống tự động lưu sample khi đủ 30 frame liên tiếp có landmarks hợp lệ.
  \item Thanh tiến độ hiển thị số sample đã thu; hoàn thành khi đạt 150 samples.
\end{enumerate}

\textit{Lý do phát nhạc nền khi thu:}
Thay vì yêu cầu người thu ``diễn xuất'' theo hướng dẫn chữ (dễ gây chuyển động
gượng gạo, thiếu tự nhiên), nhạc nền kích thích cảm xúc thật của người thu,
tạo ra chuyển động cơ thể phản ánh đúng trạng thái cảm xúc mong muốn.
Đây là một đặc điểm thiết kế quan trọng giúp nâng cao chất lượng nhãn.

\subsection{Phân Tích Thống Kê Dataset}

\begin{table}[H]
  \centering
  \caption{Thống kê tổng quan dataset cử chỉ -- cảm xúc}
  \label{tab:dataset_stats}
  \begin{tabular}{lc}
    \toprule
    \textbf{Thuộc tính} & \textbf{Giá trị} \\
    \midrule
    Số người thu thập     & 2 (A, C) \\
    Tổng số mẫu           & 1.853 \\
    Số mẫu trung bình/chế độ & $\approx 309$ \\
    Số frame mỗi mẫu      & 30 ($\approx 1$ giây ở 30 fps) \\
    Số features mỗi frame & 225 \\
    Tổng frames           & 55.590 \\
    Kích thước file CSV   & $\approx 85$ MB \\
    \bottomrule
  \end{tabular}
\end{table}

\begin{figure}[H]
  \centering
  \includegraphics[width=\textwidth]{images/data/fig02_dataset_stats.png}
  \caption{Thống kê và chất lượng dataset.
           (a)~Số lượng mẫu mỗi chế độ cho thấy phân bố tương đối đều.
           (b)~Tỉ lệ frame có landmarks hợp lệ: các chế độ có tay đạt $>85\%$,
           chế độ NONE đạt $<5\%$ (đúng theo thiết kế).
           (c)~Phân bố giá trị landmarks chuẩn hóa xung quanh 0.}
  \label{fig:dataset_stats}
\end{figure}

\begin{figure}[H]
  \centering
  \includegraphics[width=\textwidth]{images/data/fig03_time_series.png}
  \caption{Phân tích chuỗi thời gian cho từng chế độ cảm xúc.
           Mỗi biểu đồ hiển thị magnitude trung bình của landmarks bàn tay phải
           theo 30 frame (~1 giây). Đường đậm = trung bình, vùng bóng = $\pm 1$ std.
           HAPPY\_HIGH có biên độ lớn và biến động cao;
           SAD\_LOW có biên độ thấp và ổn định.}
  \label{fig:time_series}
\end{figure}

\begin{figure}[H]
  \centering
  \includegraphics[width=\textwidth]{images/data/fig04_motion_features.png}
  \caption{So sánh đặc trưng động học giữa 6 chế độ.
           (a)~Tốc độ di chuyển landmark theo frame: HAPPY\_HIGH có tốc độ
           cao nhất và biến động nhiều nhất; SAD\_LOW có tốc độ thấp nhất.
           (b)~Phân bố tốc độ trung bình mỗi sample dưới dạng boxplot,
           cho thấy sự phân tách rõ ràng giữa các chế độ.}
  \label{fig:motion}
\end{figure}

\begin{figure}[H]
  \centering
  \includegraphics[width=\textwidth]{images/data/fig05_correlation_pca.png}
  \caption{Phân tích tương quan và PCA.
           (a)~Ma trận tương quan 9 features đầu: các tọa độ cùng ngón tay
           ($x$, $y$, $z$) có tương quan cao (vùng xanh đậm).
           (b)~PCA 2 chiều trên 500 mẫu: các chế độ HAPPY\_HIGH và SAD\_LOW
           tách biệt rõ nhất trên không gian PC1-PC2.}
  \label{fig:pca}
\end{figure}

% ────────────────────────────────────────────────────────────────────
\section{Lakh MIDI Clean Dataset}
% ────────────────────────────────────────────────────────────────────

\subsection{Nguồn Gốc và Lý Do Lựa Chọn}

Music Prior cần được huấn luyện trên nhạc do con người sáng tác
để học được ngữ pháp âm nhạc thực sự, không phải các pattern nhân tạo.
\textbf{Lakh MIDI Clean Dataset}~\cite{raffel2016} là tập con đã được làm sạch
của Lakh MIDI Dataset gốc, gồm các bản nhạc MIDI chất lượng cao
từ nhiều thể loại (pop, rock, jazz, classical).

\subsection{Tiền Xử Lý}

\begin{enumerate}
  \item \textbf{Lấy mẫu:} Chọn ngẫu nhiên 5.000 bản nhạc.
  \item \textbf{Trích xuất melody:} Track có nhiều nốt nhất trong mỗi bản
        được chọn làm melody đại diện.
  \item \textbf{Token hóa:}
        \begin{itemize}
          \item Nốt MIDI $0$--$127$ $\rightarrow$ token tương ứng.
          \item Khoảng lặng $> 0.2$s $\rightarrow$ token REST (128).
          \item Padding $\rightarrow$ token PAD (129).
        \end{itemize}
  \item \textbf{Sliding window:} Cắt mỗi melody thành đoạn 128 token,
        bước nhảy 64 token.
\end{enumerate}

\begin{table}[H]
  \centering
  \caption{Thống kê Lakh MIDI sau tiền xử lý}
  \label{tab:midi_stats}
  \begin{tabular}{lc}
    \toprule
    \textbf{Thuộc tính} & \textbf{Giá trị} \\
    \midrule
    Số bản nhạc sử dụng       & 5.000 \\
    Vocabulary size            & 130 (128 nốt + REST + PAD) \\
    Chiều dài mỗi sequence    & 128 tokens \\
    Tỉ lệ train/val            & 90\% / 10\% \\
    \bottomrule
  \end{tabular}
\end{table}
